%% ch04 --- Objects and data: dunder methods, equality, hashing, records
%%
%% Written September 2026. Every listing under code/ch04/ runs on the pinned
%% interpreter and is executed by code/tests/test_listings.py; every quoted
%% output is a generated transcript. This chapter owes no measurement: every
%% claim in it is about semantics, and is checked by a test rather than timed.

\chapter{Objects and data: dunder methods, equality, hashing, records}
\label{ch:objects}

\noindent
C\# gives you a named feature for each thing an object can do. An indexer
is a language construct; \code{Equals} and \code{GetHashCode} are virtual
methods you override in pairs; \code{IComparable} is an interface you
declare; \code{record} is a keyword that writes several of these for you.
Python has all of the same capabilities and none of the named features. It
has one mechanism instead: a method with two underscores at each end, which
the syntax calls on your behalf.

The mechanism is worth one sentence before the detail. \textbf{The protocol
is the method, and nothing declares it.} There is no interface to
implement, no base class to inherit and no attribute to apply: a class that
defines \code{\_\_lt\_\_} can be sorted, a class that defines
\code{\_\_len\_\_} answers \code{len()}, and a class that defines neither is
simply never asked. That is the payoff of this chapter \dash{} you can read
any library's class, look at the dunder methods on it, and name the C\#
feature each one is imitating.

It is also where the chapter's traps come from. A duck-typed protocol
cannot warn you that you have half-implemented it, so the compiler that
would have told you is not there.

\section{Every operator is a method}
\label{sec:ch04-protocol}

\index{dunder methods}
The mapping is not this book's invention and does not have to be taken on
trust: the standard library writes it down. Every infix operator has a
function in \api{operator}, and every one of those functions is spelled as
a method on the left-hand operand.

\pyregion{code/ch04/dunders.py}{protocol}{The operator, the method, and a check that the two agree}{lst:ch04-protocol}

\mermaidfig{ch04-dunder-map}{A named C\# feature becomes ordinary syntax, which calls a method that may or may not be there}{how a C\# feature arrives in Python}

\noindent
The last column of that listing is the point. Nothing asserts the
equivalence; each row runs the operator function and the method it names
and reports that the two agreed.

\index{\_\_repr\_\_@\texttt{\_\_repr\_\_}}
\index{\_\_str\_\_@\texttt{\_\_str\_\_}}
Two of these look like one C\# method and are not. \code{ToString()} is a
single override that every caller gets; Python splits it, and splits it the
way round that catches people. \code{\_\_str\_\_} is what \code{print} and
\code{format} show. \code{\_\_repr\_\_} is what the REPL shows, what a
debugger shows, and \dash{} the half that does the damage \dash{} what a
\emph{list} of your objects shows, because a container formats its elements
with \code{repr} rather than \code{str}.

\pyregion{code/ch04/dunders.py}{repr}{One of these is the one you owe}{lst:ch04-repr}

\transcript{ch04-dunders}

\begin{trapbox}
\textbf{\code{\_\_str\_\_} is not \code{ToString()}; \code{\_\_repr\_\_} is
closer.} A class with only \code{\_\_str\_\_} prints beautifully on its own
and prints as \code{<module.Class object at 0x...>} the moment it is inside
a list, a dict, a tuple or a log record \dash{} which is where you were
going to read it. Write \code{\_\_repr\_\_} first. \code{\_\_str\_\_} is
optional and falls back to it; the reverse is not true.
\end{trapbox}

\section{The method that disappears}
\label{sec:ch04-equality}

\index{equality!and hashing}
This is the chapter's headline, and the one mapping where C\# is genuinely
more forgiving. Override \code{Equals} there and not \code{GetHashCode} and
the compiler raises CS0659 \dash{} a warning, which you have read many
times and possibly ignored many times, because the type still works. The
dictionary degrades and nothing breaks.

Python does not warn. It acts.

\pyregion{code/ch04/equality.py}{broken}{Value equality, and nothing else. This class cannot be a dict key}{lst:ch04-broken}

\noindent
Nothing in that class mentions hashing. Defining \code{\_\_eq\_\_} is what
removed it: Python sets \code{\_\_hash\_\_} to \code{None} on any class that
defines \code{\_\_eq\_\_} without defining \code{\_\_hash\_\_}, on the
reasoning that an object whose equality you have redefined can no longer
honour the identity hash it inherited.

\mermaidfig{ch04-eq-hash}{Defining one method removes another, and the failure arrives somewhere else}{what defining \code{\_\_eq\_\_} costs}

\transcript{ch04-equality}

\begin{trapbox}
\textbf{Defining \code{\_\_eq\_\_} makes your class unhashable.} Not
slower, not subtly wrong: \code{\_\_hash\_\_} becomes \code{None} and the
object cannot be a \code{dict} key or a \code{set} member at all. The
failure is a \code{TypeError} at the first container that touches it, which
is usually a long way from the class. The fix is to write both, over one
key tuple, so they cannot disagree.
\end{trapbox}

\pyregion{code/ch04/equality.py}{fixed}{Both methods, over one key tuple, so they cannot drift apart}{lst:ch04-fixed}

\begin{csbox}
The contract is the one you already know and it has not changed: equal
objects must hash equal, and an object's hash must not change while it is
in a dictionary. What has changed is the enforcement. C\# warns and
continues; Python removes the capability. And there is one thing Python
asks for that C\# does not \dash{} returning \code{NotImplemented} rather
than \code{False} when the other operand is of a type you do not know,
which lets the interpreter try the other object's \code{\_\_eq\_\_} before
deciding the two are unequal.
\end{csbox}

\noindent
\index{is@\code{is} against \code{==}}
The neighbouring habit is \code{is}, and the folklore about it is worth
correcting because it is repeated everywhere. \code{is} asks whether two
names refer to the same object, and never asks anything else. The usual
warning is that it happens to work for small integers and breaks one
magnitude up. Run that demonstration inside a function on the pinned
interpreter and it does not reproduce: \code{a = 257; b = 257; a is b} is
\code{True}.

\pyregion{code/ch04/equality.py}{identity}{Two mechanisms, and only one of them is the cache everybody names}{lst:ch04-identity}

\noindent
Two separate things are going on, and neither is the one the folklore
names. CPython caches the integers \code{-5} to \code{256}, so a value
built at run time really is the same object as a literal in that range.
Independently of that, the compiler stores equal constants in a code
object's constant table \emph{once}, so two literal \code{257}s in one
function body are one object \dash{} which is why the classic
demonstration works in a REPL, where each statement is its own code object,
and fails in a script.

\begin{trapbox}
\textbf{\code{is} on integers hides from the test you would write for it.}
The bug needs a value that crossed a run-time boundary, and that is where
every real value comes from: a parsed request, a database row, a JSON
payload. So \code{is} looks correct in every small example and fails in
production. Keep \code{is} for \code{None} and for sentinels, where
identity is the question you actually mean.
\end{trapbox}

\begin{exercise}{e04_01_hashable}{Value equality that survives a set}
Write \code{\_\_eq\_\_} and \code{\_\_hash\_\_} on a \code{Seat} so that
equal seats compare equal, are usable as \code{dict} keys, and collapse to
one entry in a \code{set}. Five tests decide it, and the one that catches a
half-finished answer is the last: an implementation that stops after
\code{\_\_eq\_\_} passes the first test and fails everything after it.
\end{exercise}

\section{A record, and what a dataclass gives you}
\label{sec:ch04-dataclass}

\index{dataclass}
A C\# \code{record} is a keyword that writes value equality, a hash, a
printable form and a \code{with} expression. \code{@dataclass} is the same
idea as a decorator, and reading the hand-written version beside it is the
quickest way to see exactly which methods a record is.

\pyregion{code/ch04/money.py}{byhand}{Every method here is one the decorator would have written}{lst:ch04-byhand}

\pyregion{code/ch04/money.py}{dataclass}{The same type. The flags are the interesting part}{lst:ch04-dataclass}

\transcript{ch04-money}

\noindent
\code{frozen=True} is the half that makes it a value, and it is also what
permits \code{@dataclass} to write \code{\_\_hash\_\_} at all, for the
reason the previous section gave.
\code{FrozenInstanceError} is what an assignment raises instead. \api{dataclasses.replace} is \code{with}: a new object rather than a
mutation, and it runs \code{\_\_init\_\_}, so validation is not skipped and
a field added later is carried without the method changing.

\code{\_\_slots\_\_} has no C\# feature to map onto. It replaces the
per-instance dictionary with a fixed layout, which is the nearest thing to
what a struct costs \dash{} and it has a second effect worth more than the
memory: a typo in an attribute name raises \code{AttributeError} instead of
quietly creating a new field.

\begin{csbox}
Two things a record gives you that \code{@dataclass} does not, unless
asked. Ordering is \code{order=True}, which writes the four comparison
methods from the fields in declaration order, the way a tuple compares; in
C\# you implement \code{IComparable} by hand, so this is the rare case
where Python asks for less. And a record's \code{with} is compile-time
checked against the field names, where \api{dataclasses.replace} raises at
run time \dash{} though it does raise, which \code{setattr} on a plain
class would not.
\end{csbox}

\begin{exercise}{e04_02_reading}{A frozen record, ordered, replaced}
Declare \code{Reading} so that it is immutable, hashable, sorts by its
first field and then its second with no key function, and grows a
\code{warmer\_by} that returns a new reading. Most of the answer is the
decorator's arguments; the rest is the order the two fields are declared
in.
\end{exercise}

\section{Properties, indexers and \texorpdfstring{\code{IEnumerable}}{IEnumerable}}
\label{sec:ch04-properties}

\index{property@\code{@property}}
A C\# property is a pair of methods behind a field-like name.
\code{@property} is the same trick, with one difference in the advice that
follows from it. In C\# a field and a property differ in the compiled
interface, so widening a field later is a breaking change and you are told
to start with a property you do not yet need. In Python the attribute and
the property are spelled identically at every call site, so you start with
a plain attribute and promote it when a reason appears.

\pyregion{code/ch04/protocols.py}{property}{\code{months} began as a plain attribute; no caller changed}{lst:ch04-property}

\noindent
An indexer is \code{\_\_getitem\_\_}. \code{Count} is \code{\_\_len\_\_}.
\code{IEnumerable<T>} is \code{\_\_iter\_\_}. None of the three is declared
anywhere, and a class may implement any subset: the syntax looks for its
own method and uses it if it is there.

\pyregion{code/ch04/protocols.py}{container}{Four protocols, no interface, no base class}{lst:ch04-container}

\begin{note}
\code{\_\_contains\_\_} is the optional one, and optional in an interesting
way: leave it out and \code{in} still works, by walking
\code{\_\_iter\_\_} and comparing each element. So a class can answer
\code{in} correctly and slowly without anybody writing the method \dash{}
and a class that promises something other than exact comparison, such as
case-insensitive membership, must write it, or the promise is quietly
false.
\end{note}

\begin{exercise}{e04_03_roster}{Be a collection without declaring one}
Give \code{Roster} the methods behind \code{len()}, indexing, iteration and
\code{in}, plus a read-only computed \code{first}. Two of the five tests
are the ones that matter: membership has to ignore case, which is the only
method not merely delegating, and \code{first} has to be a property, which
the test checks by trying to assign to it.
\end{exercise}

\section{An enum that is not an integer}
\label{sec:ch04-enums}

\index{Enum@\code{Enum}}
A C\# enum is a named integer and behaves like one everywhere: it compares
to an integer, it serialises as one, and a value nobody declared is still
assignable to it. Python's \code{Enum} is a class whose members are
singletons, so a member is \emph{not} its value.

\pyregion{code/ch04/enums.py}{kinds}{Three kinds, and only one of them is a string}{lst:ch04-enums}

\transcript{ch04-enums}

\noindent
\code{Severity.LOW == 1} is \code{False}, and the checker Chapter~\ref{ch:packaging} put
in the toolchain says so before the program runs: pyright rejects the
comparison as one whose types have no overlap. That is the C\# habit caught
at the moment of typing rather than in production.

\code{StrEnum} members \emph{are} strings, which is the one to reach for
when the value crosses a wire, because \api{json.dumps} serialises it
without being asked. A plain \code{Enum} member raises \code{TypeError}
there and needs \code{.value} \dash{} the line everyone forgets, because in
C\# the enum went out as a number and nobody had to ask. \code{Flag} is
\code{[Flags]}, with \code{auto()} handing out the powers of two that get
written by hand and wrong once.

\section{Where the state lives}
\label{sec:ch04-state}

\index{class attributes}
A class attribute and a \code{static} field are both shared, so it is
tempting to call them the same thing. The difference is not the sharing. It
is that a class attribute is readable \emph{through the instance}, with no
qualification and no warning, so \code{self.seen} finds the object on the
class \dash{} and mutating it mutates it for every instance that will ever
exist.

\pyregion{code/ch04/state.py}{shared}{One list, every instance, and nothing at the call site says so}{lst:ch04-shared}

\mermaidfig{ch04-attribute-lookup}{The lookup that makes a class attribute look like a field}{how \code{self.seen} finds the class's object}

\transcript{ch04-state}

\begin{trapbox}
\textbf{A mutable class attribute is shared by accident, through
\code{self}.} The attribute lookup walks the instance's own dictionary
first, finds nothing, walks the class's, finds the list, and mutates that
one. In C\# you would have had to write the type's name to reach a static,
and that is exactly the tell Python does not give you. Per-instance state
is created in \code{\_\_init\_\_} and nowhere else.
\end{trapbox}

\noindent
\code{@dataclass} is the one place in the language that stops this before
it runs: a mutable default raises \code{ValueError} at class creation,
naming \code{default\_factory}. It cannot see the plain-class spelling
above, which is the spelling this trap actually arrives in.

The same aliasing appears without a class in sight. \code{[[]] * n} does
not give you \code{n} lists; it gives one list \code{n} times, because
\code{*} copies the reference. A comprehension evaluates its body once per
item and gives you \code{n}.

\pyregion{code/ch04/state.py}{aliasing}{The multiplication copies a reference; the comprehension does not}{lst:ch04-aliasing}

\begin{exercise}{e04_04_shared}{Shared by accident, shared on purpose}
\code{Basket} keeps its items in a class attribute, so every basket shares
one list. Fix it, and then keep a counter of how many baskets have been
opened \dash{} which is shared \emph{deliberately} and is what a
\code{static} field really maps to. The second half is the one worth
slowing down for: incrementing it through \code{self} would read the
class's value and bind a new instance attribute, leaving the class's count
at zero for ever.
\end{exercise}

\section{Two operators that are not C\#'s}
\label{sec:ch04-arithmetic}

\index{division!floor against truncating}
\code{/} on two integers is integer division in C\# and true division in
Python, which is the difference everybody is told about. The one that is
not mentioned is that \code{//} is not C\#'s \code{/} either. It
\emph{floors}, towards negative infinity, where C\# truncates towards zero
\dash{} so the two languages disagree about $-7$ divided by $2$, and no
cast repairs it. C\#'s answer is Python's \code{int(a / b)}.

\pyregion{code/ch04/arithmetic.py}{division}{Four answers where C\# gives two}{lst:ch04-division}

\transcript{ch04-arithmetic}

\noindent
The same split runs through the remainder: Python's \code{\%} takes the
sign of the divisor, matching \code{//}, and \api{math.fmod} takes the sign
of the dividend, matching C\#'s \code{\%}. Both are available; only one is
the operator.

\begin{trapbox}
\textbf{\code{0.1 + 0.2 != 0.3} is not a Python defect.} It is IEEE 754,
and it is equally true in C\#. The printing used to differ and no longer
does: .NET Framework formatted a \code{double} to fifteen significant
digits and hid the error, and since .NET Core 3.0 the default has been the
shortest round-tripping string \dash{} which is what Python has always
printed. Both now show \code{0.30000000000000004}. Neither language is more
accurate than the other. Use \api{math.isclose} for computed floats, and
\api{decimal.Decimal} for money, which is the same advice as C\#'s
\code{decimal}.
\end{trapbox}

\section{What this bought you}
\label{sec:ch04-summary}

\begin{center}
\small
\begin{tabularx}{\linewidth}{@{}lX@{}}
\toprule
\textbf{C\#} & \textbf{Python} \\
\midrule
operator overloading & \code{\_\_add\_\_}, \code{\_\_lt\_\_} and the rest,
  none declared \\
\code{ToString()} & \code{\_\_repr\_\_}, and optionally
  \code{\_\_str\_\_} \\
\code{Equals} + \code{GetHashCode} & \code{\_\_eq\_\_} +
  \code{\_\_hash\_\_}, and omitting the second removes it \\
\code{IComparable} & \code{\_\_lt\_\_}, or \code{order=True} \\
\code{record}, \code{with} & \code{@dataclass(frozen=True)},
  \api{dataclasses.replace} \\
a property & \code{@property}, promoted from a plain attribute \\
an indexer, \code{Count} & \code{\_\_getitem\_\_}, \code{\_\_len\_\_} \\
\code{IEnumerable<T>} & \code{\_\_iter\_\_} \\
\code{enum}, \code{[Flags]} & \code{Enum}, \code{StrEnum}, \code{Flag} \\
a \code{static} field & a class attribute \dash{} but readable through
  \code{self} \\
\code{/} on two \code{int}s & \code{//}, which floors rather than
  truncating \\
\bottomrule
\end{tabularx}
\end{center}

\noindent
The payoff is the one the chapter opened with, and it is a reading skill
rather than a writing one. Open any class in any library, look at the
methods with underscores at both ends, and you can say what it is for: this
one is a value type, that one is a collection, this one is a context
manager. Nothing in the class says so, because nothing has to. The price of
that economy is the section that follows the mechanism in each case
\dash{} a protocol nothing declares is a protocol nothing can tell you that
you have half-finished, and the two halves of \code{\_\_eq\_\_} and
\code{\_\_hash\_\_} are where that bill arrives first.
